\chapter{A Short History of Secret Writing}

In Chapter~1, we saw that modern cryptography demands precise definitions of secrecy and computationally bounded adversaries. Before building the mathematics of those modern guarantees, we need to understand where they came from -- and why intuitive attempts to hide secrets failed. For most of human history, hiding a message meant moving letters around. There were no
computers, no notion of a ``key space'' measured in bits, and no proofs. There were clever
people on one side trying to keep a secret, and equally clever people on the other side
trying to read it. That contest is worth studying closely. Every cipher in this
chapter was, in its day, considered strong by serious people, and every one of them was
eventually torn open. The way each one fell teaches a lesson that survives into the modern
schemes we build later in this book: where the secret really lives, why reusing it is fatal,
and why ``it looks scrambled to me'' is not a security argument.

We will work the classic breaks by hand, with small numbers and real letters, because the
intuition behind frequency analysis and the Kasiski test is exactly the same intuition you
need to reason about modern attacks. The tools change; the way an attacker thinks does not.

\chapterline

\section{The cast and the rules of the game}
\label{ch02:cast}

Recall from Chapter~1 our cast of characters and the four standard threat models. \textbf{Alice} wants to send a message to \textbf{Bob}; \textbf{Eve} sits on the wire between them and copies down everything that passes. Alice and Bob share some secret in advance; Eve does not have it. The question is always the same: can Eve recover the message anyway?

It helps to view classical ciphers in terms of these threat models, because almost every historical failure came from underestimating Eve's access.

\begin{definition}[Attack models in historical context]
\label{ch02:attack-models}
In Chapter~1, we classified an attacker by what information she commands. In historical cryptanalysis, these models emerge in direct sequence:
\begin{itemize}
  \item \emph{Ciphertext-only.} Eve sees only encrypted messages. This was the default position for ancient military couriers and classical frequency analysis.
  \item \emph{Known-plaintext.} Eve sees some plaintext/ciphertext pairs (for example, she guesses that a message ends ``Heil Hitler'' or begins ``WEATHER REPORT'').
  \item \emph{Chosen-plaintext.} Eve can get messages of her choosing encrypted under the unknown key.
  \item \emph{Chosen-ciphertext.} Eve can get ciphertexts of her choosing decrypted (or at least learn whether they decrypt validly).
\end{itemize}
\end{definition}

\begin{focusbox}
A modern scheme must survive chosen-plaintext (and usually chosen-ciphertext) attacks.
Every classical cipher in this chapter dies to a ciphertext-only attack -- the weakest model
there is. Keep that gap in mind; it measures how far the field traveled.
\end{focusbox}

There is one more principle worth keeping in mind, because the entire history of secret writing validates it: Kerckhoffs's principle.

\begin{formalbox}
\textbf{Kerckhoffs's principle (1883).} A cipher should remain secure even if everything
about the system except the key is public knowledge. The Dutch linguist Auguste Kerckhoffs
laid out six rules for military ciphers in the \emph{Journal des sciences militaires}; this
is the one that endured. Claude Shannon restated it crisply: ``the enemy knows the system.''
The practical reading: do not bet your secret on the secrecy of your \emph{method}. Methods
leak, get captured, get reverse-engineered, or get sold. Only the key should be secret, and
the key should be easy to change. Nearly every disaster in this chapter is a violation of
this rule.
\end{formalbox}

\section{The Caesar cipher: shift and pray}
\label{ch02:caesar}

The oldest cipher most people know is named for Julius Caesar, who according to Suetonius
shifted each letter of his military dispatches three places down the alphabet. \texttt{A}
becomes \texttt{D}, \texttt{B} becomes \texttt{E}, and so on, wrapping \texttt{X}, \texttt{Y},
\texttt{Z} around to \texttt{A}, \texttt{B}, \texttt{C}.

To do anything with this mathematically, number the letters $\texttt{A}=0, \texttt{B}=1,
\ldots, \texttt{Z}=25$. Encryption with shift $k$ is then arithmetic modulo $26$:
\[
  c \equiv m + k \pmod{26}, \qquad m \equiv c - k \pmod{26}.
\]
Caesar used $k=3$. The key is just the single number $k$, and there are only $26$ possible
shifts (one of which, $k=0$, does nothing).

\begin{example}[Encrypting by hand]
\label{ch02:caesar-example}
Take the plaintext \texttt{ATTACK} and $k=3$.
\[
\begin{array}{c|cccccc}
\text{plain}  & \texttt{A} & \texttt{T} & \texttt{T} & \texttt{A} & \texttt{C} & \texttt{K}\\
m             & 0 & 19 & 19 & 0 & 2 & 10\\
m+3 \bmod 26  & 3 & 22 & 22 & 3 & 5 & 13\\
\text{cipher} & \texttt{D} & \texttt{W} & \texttt{W} & \texttt{D} & \texttt{F} & \texttt{N}
\end{array}
\]
So \texttt{ATTACK} encrypts to \texttt{DWWDFN}. To decrypt, subtract $3$: \texttt{D}$\to$%
\texttt{A}, \texttt{W}$\to$\texttt{T}, and you are back.
\end{example}

A shift cipher is easiest to picture as two concentric alphabet rings: line up the plaintext
\texttt{A} on the inner ring with the ciphertext letter $k$ positions around on the outer ring,
and reading off any plaintext letter gives its ciphertext. Figure~\ref{ch02:fig-caesar-wheel}
shows the $k=3$ alignment of Example~\ref{ch02:caesar-example}.

\begin{figure}[H]
\centering
\begin{tikzpicture}[font=\footnotesize,>=stealth,scale=0.95]
  \def\Rout{2.65}\def\Rmid{1.90}\def\Rin{1.15}
  
  % Circular bands
  \draw[thick,quotientorange!50,fill=quotientorange!8] (0,0) circle (\Rout);
  \draw[thick,sourceblue!50,fill=sourceblue!8] (0,0) circle (\Rmid);
  \draw[thick,structurepurple!40,fill=white] (0,0) circle (\Rin);

  % Inner ring = plaintext A..Z ; outer ring = ciphertext (i+3) mod 26
  \foreach \i/\pl in
    {0/A,1/B,2/C,3/D,4/E,5/F,6/G,7/H,8/I,9/J,10/K,11/L,12/M,
     13/N,14/O,15/P,16/Q,17/R,18/S,19/T,20/U,21/V,22/W,23/X,24/Y,25/Z} {
    \pgfmathsetmacro\ang{90 - \i*360/26}
    \node[font=\bfseries\tiny,sourceblue!90!black] at (\ang:\Rin+0.36) {\texttt{\pl}};
  }
  \foreach \i/\cl in
    {0/D,1/E,2/F,3/G,4/H,5/I,6/J,7/K,8/L,9/M,10/N,11/O,12/P,
     13/Q,14/R,15/S,16/T,17/U,18/V,19/W,20/X,21/Y,22/Z,23/A,24/B,25/C} {
    \pgfmathsetmacro\ang{90 - \i*360/26}
    \node[font=\bfseries\tiny,quotientorange!90!black] at (\ang:\Rmid+0.36) {\texttt{\cl}};
  }

  % Center annotations (well separated from rotation arc)
  \node[font=\bfseries\scriptsize,structurepurple] at (0,0.28) {Shift $k=3$};
  \draw[->,thick,structurepurple] (-0.32,-0.18) arc (190:350:0.32);

  % Side explanatory box with generous line spacing and clean placement
  \node[draw=sourceblue!40,fill=sourceblue!5,rounded corners=4pt,thick,align=left,font=\scriptsize,inner sep=7pt] at (6.2,0) {
    \textbf{Concentric Disk Mechanics}\\[4pt]
    -- Inner wheel: Plaintext $m$\\[2pt]
    -- Outer wheel: Ciphertext $c \equiv m+3 \pmod{26}$\\[2pt]
    -- Keyspace: Exactly $26$ rotations\\[2pt]
    -- Brute force cost: $25$ tries
  };
\end{tikzpicture}
\caption{The Caesar cipher as two concentric alphabet rings. With $k=3$, the inner (plaintext) \texttt{A} lines up against outer (ciphertext) \texttt{D}. Because only $26$ distinct relative rotations exist, an adversary can exhaust every candidate key in seconds.}
\label{ch02:fig-caesar-wheel}
\end{figure}

\subsection{Why it fails: there is nowhere for the secret to hide}

The Caesar cipher has a key space of size $26$. An attacker simply tries all of them.

\begin{attackbox}
\textbf{Brute force in 25 tries.} Eve intercepts \texttt{WKHWUHDVXUHLVEXULHG}. She writes out
all $25$ nontrivial shifts of the first few letters and reads down the list for English:
\[
\begin{array}{rl}
k=1: & \texttt{VJGVTGCUWTG...}\\
k=2: & \texttt{UIFUSFBTVSF...}\\
k=3: & \texttt{THETREASURE...}\quad\leftarrow\text{readable!}
\end{array}
\]
Shift by $k=3$ and the message reads \texttt{THETREASUREISBURIED}. She did not need any
insight into Alice's method; she only needed to try $25$ shifts.
\end{attackbox}

The lesson recurs throughout cryptography: \emph{the number of possible keys must
be too large to enumerate.} Twenty-six is far too few. Even by hand a person clears all $25$
shifts in under a minute. This is the first appearance of a quantity we will obsess over for
the rest of the book -- the size of the key space -- and the first demonstration that if it is
small, nothing else you do matters.

\begin{connectionbox}
``Too large to enumerate'' is a statement about counting, which is the subject of the
companion combinatorics chapters. A modern $128$-bit key has $2^{128}\approx 3.4\times10^{38}$
possibilities. If you could test a trillion keys per second on a billion machines, you would
still need on the order of $10^{10}$ years -- comfortably longer than the age of the universe.
That gap, from $26$ keys to $2^{128}$, is what makes brute force hopeless against a modern
cipher. We make this computational security model precise in Chapter~8.
\end{connectionbox}

\section{General substitution: a huge key that still loses}
\label{ch02:substitution}

The obvious fix for a tiny key space is to make it enormous. Instead of shifting the whole
alphabet by one fixed amount, pick an arbitrary permutation: a secret one-to-one map from
each plaintext letter to a ciphertext letter. This is the \emph{monoalphabetic substitution
cipher}.

A worked key might be:
\[
\begin{array}{c|cccccccccccccccccccccccccc}
\text{plain}  & \texttt{A}&\texttt{B}&\texttt{C}&\texttt{D}&\texttt{E}&\texttt{F}&\texttt{G}&\texttt{H}&\texttt{I}&\texttt{J}&\texttt{K}&\texttt{L}&\texttt{M}&\texttt{N}&\texttt{O}&\texttt{P}&\texttt{Q}&\texttt{R}&\texttt{S}&\texttt{T}&\texttt{U}&\texttt{V}&\texttt{W}&\texttt{X}&\texttt{Y}&\texttt{Z}\\
\text{cipher} & \texttt{Q}&\texttt{W}&\texttt{E}&\texttt{R}&\texttt{T}&\texttt{Y}&\texttt{U}&\texttt{I}&\texttt{O}&\texttt{P}&\texttt{A}&\texttt{S}&\texttt{D}&\texttt{F}&\texttt{G}&\texttt{H}&\texttt{J}&\texttt{K}&\texttt{L}&\texttt{Z}&\texttt{X}&\texttt{C}&\texttt{V}&\texttt{B}&\texttt{N}&\texttt{M}
\end{array}
\]
Now the key is the entire bottom row, and there are
\[
  26! = 403{,}291{,}461{,}126{,}605{,}635{,}584{,}000{,}000 \approx 4 \times 10^{26}
\]
possible keys. That is about $88$ bits. Brute force is hopeless -- you cannot try
$4\times10^{26}$ alphabets by hand or, in 1500, by any means at all. For centuries people
believed this made the cipher unbreakable. They were wrong, and the reason is the most
important idea in classical cryptanalysis.

\subsection{Frequency analysis: the key space is irrelevant if the cipher leaks}

Here is the flaw. A monoalphabetic substitution maps every \texttt{E} in the plaintext to the
same ciphertext letter, every \texttt{T} to the same letter, and so on. It scrambles
\emph{which} symbol stands for \texttt{E}, but it cannot hide \emph{how often} that symbol
appears. English is lopsided: \texttt{E} is about $12.7\%$ of letters, \texttt{T} about
$9.1\%$, \texttt{A} about $8.2\%$, while \texttt{J}, \texttt{Q}, \texttt{X}, \texttt{Z} are
each well under $1\%$. The cipher carries this fingerprint straight through.

\begin{intuitionbox}
Imagine a room where everyone swaps name tags while nobody changes
height. If you know the tallest person in the room is Eve, you find Eve by looking for the
tallest person, no matter what her name tag says. Letter frequencies are heights the substitution
cannot change. Count the most common ciphertext symbol; it is overwhelmingly likely to be
\texttt{E}. The 88-bit key did not help, because we never had to guess it -- we read it off
the statistics.
\end{intuitionbox}

\paragraph{Al-Kindi and the birth of frequency analysis (c.~850 CE).} The earliest surviving description of frequency
analysis is by the Arab polymath Abu Yusuf al-Kindi, in a manuscript on deciphering encrypted
messages written around 850 CE. His insight: count the symbols, line them up against the known
frequencies of the language, and the substitution unravels. This was the
first general cipher-breaking technique in history, and it predates the European cryptography
it would later demolish by six centuries. It was also the first time anyone used \emph{statistics}
as a weapon, an idea that runs straight through to modern cryptanalysis.

\begin{figure}[H]
\centering
\begin{tikzpicture}[font=\small,scale=0.95]
  % Y axis
  \draw[->,thick,gray!70] (0,0) -- (0,3.2) node[above,font=\tiny\bfseries]{Freq (\%)};
  \foreach \pct/\y in {0/0, 5/1.1, 10/2.2} {
    \draw[gray!60] (-0.1,\y) -- (0,\y);
    \node[left,font=\tiny,gray!80] at (-0.1,\y) {\pct\%};
  }

  % Uniform reference line
  \draw[dashed,thick,trapred!80] (0,0.847) -- (10.6,0.847)
    node[right,font=\tiny\bfseries,trapred!90!black]{Uniform noise $1/26 \approx 3.85\%$};

  % Bars for letters A-Z (x = 0.2 + i * 0.39)
  % Frequencies: A=8.2, B=1.5, C=2.8, D=4.3, E=12.7, F=2.2, G=2.0, H=6.1, I=7.0,
  % J=0.2, K=0.8, L=4.0, M=2.4, N=6.7, O=7.5, P=1.9, Q=0.1, R=6.0, S=6.3, T=9.1,
  % U=2.8, V=1.0, W=2.4, X=0.2, Y=2.0, Z=0.1
  \foreach \i/\L/\h/\cat in {
    0/A/1.80/1, 1/B/0.33/0, 2/C/0.62/0, 3/D/0.95/0,
    4/E/2.79/2, 5/F/0.48/0, 6/G/0.44/0, 7/H/1.34/1,
    8/I/1.54/1, 9/J/0.05/0, 10/K/0.18/0, 11/L/0.88/0,
    12/M/0.53/0, 13/N/1.47/1, 14/O/1.65/1, 15/P/0.42/0,
    16/Q/0.02/0, 17/R/1.32/1, 18/S/1.39/1, 19/T/2.00/3,
    20/U/0.62/0, 21/V/0.22/0, 22/W/0.53/0, 23/X/0.04/0,
    24/Y/0.44/0, 25/Z/0.02/0
  } {
    \pgfmathsetmacro\xpos{0.2 + \i * 0.39}
    \ifcase\cat
      \draw[fill=gray!30,draw=gray!80!black,rounded corners=0.5pt] (\xpos,0) rectangle (\xpos+0.28,\h);
    \or
      \draw[fill=sourceblue!35,draw=sourceblue!90!black,rounded corners=0.5pt] (\xpos,0) rectangle (\xpos+0.28,\h);
    \or
      \draw[fill=quotientorange!50,draw=quotientorange!95!black,rounded corners=0.5pt] (\xpos,0) rectangle (\xpos+0.28,\h);
    \or
      \draw[fill=targetgreen!50,draw=targetgreen!95!black,rounded corners=0.5pt] (\xpos,0) rectangle (\xpos+0.28,\h);
    \fi
    \node[below,font=\tiny\bfseries] at (\xpos+0.14,-0.05) {\texttt{\L}};
  }

  % Annotations pointing out E and T peaks
  \node[above,font=\tiny\bfseries,quotientorange!90!black] at (0.2 + 4*0.39 + 0.14, 2.82) {12.7\%};
  \node[above,font=\tiny\bfseries,targetgreen!90!black] at (0.2 + 19*0.39 + 0.14, 2.03) {9.1\%};

  % Bottom label
  \draw[thick,gray!70] (0,0) -- (10.6,0);
\end{tikzpicture}
\caption{The frequency signature of written English. Monoalphabetic substitution merely permutes these bar heights across different letters, preserving the overall shape of the distribution. An attacker identifies high peaks (such as \texttt{E} and \texttt{T}) and rare letters (\texttt{Q}, \texttt{X}, \texttt{Z}) to recover the key without brute force.}
\label{ch02:fig-freq-hist}
\end{figure}

\subsection{Breaking one by hand}
\label{ch02:break-substitution}

Let us actually do it. Eve intercepts the following ciphertext (spaces left in only so you can
follow along; a real interceptor would strip them, and the attack still works):

\begin{quote}\ttfamily
ZG WT GK FGZ ZG WT ZIQZ OL ZIT JXTLZOGF
\end{quote}

\textbf{Step 1: count letters.} Tally every symbol (ignore spaces; there are $30$ letters).
\[
\begin{array}{l}
\texttt{Z}: 7 \quad \texttt{G}: 5 \quad \texttt{T}: 4 \quad \texttt{W}: 2 \quad \texttt{F}: 2
\quad \texttt{I}: 2 \quad \texttt{O}: 2 \quad \texttt{L}: 2\\
\texttt{K}: 1 \quad \texttt{Q}: 1 \quad \texttt{J}: 1 \quad \texttt{X}: 1
\end{array}
\]
\texttt{Z} dominates with $7$ occurrences, then \texttt{G} with $5$ and \texttt{T} with $4$.
In English, the runaway leaders are \texttt{E} and \texttt{T}, with \texttt{O} and \texttt{A}
close behind.

\textbf{Step 2: exploit structure, not just counts.} Raw frequency is a starting hypothesis,
not the whole attack. The real power comes from short words and repeated patterns. Notice the
two-letter word \texttt{ZG} appears twice, and a three-letter word \texttt{ZIT} appears. The
most common English two-letter words are \texttt{TO}, \texttt{OF}, \texttt{IN}, \texttt{IS},
\texttt{IT}, \texttt{BE}; the most common three-letter word is \texttt{THE}. A three-letter word
whose first letter is the most common ciphertext letter is a prime suspect for \texttt{THE}.
\texttt{ZIT} fits if $\texttt{Z}\to\texttt{T}$, $\texttt{I}\to\texttt{H}$, $\texttt{T}\to%
\texttt{E}$. That is consistent with the counts: \texttt{Z} was the most common symbol (so
$\texttt{Z}\to\texttt{T}$ is plausible) and \texttt{T} was third (so $\texttt{T}\to\texttt{E}$
fits a high-frequency letter). It also reads the repeated two-letter word as
$\texttt{ZG}=\texttt{T?}$, and the only common English two-letter word starting with \texttt{T}
is \texttt{TO}, forcing $\texttt{G}\to\texttt{O}$ -- and \texttt{G} was the second most common
symbol, exactly where \texttt{O} sits in English. Three independent checks agree. Promising.

\textbf{Step 3: propagate guesses.} Assume $\texttt{Z}\to\texttt{t}$, $\texttt{I}\to\texttt{h}$,
$\texttt{T}\to\texttt{e}$, $\texttt{G}\to\texttt{o}$. Rewrite, lowercase for known plaintext:
\begin{quote}\ttfamily
to We oK Fot to We tIQt OL the JXeLtOoF
\end{quote}
The repeated word \texttt{to We} now appears twice and \texttt{Fot} sits between them.
\texttt{Fot} is almost certainly \texttt{not}, giving $\texttt{F}\to\texttt{n}$. Then \texttt{We}
is \texttt{?e}; after ``to be or not to be'' leaps to mind, try $\texttt{W}\to\texttt{b}$, so
\texttt{We}$=$\texttt{be}, and \texttt{oK} becomes \texttt{o?} $=$ \texttt{or} with
$\texttt{K}\to\texttt{r}$. Rewrite:
\begin{quote}\ttfamily
to be or not to be tIQt OL the JXeLtOoF
\end{quote}
Now \texttt{tIQt} is \texttt{t?et}-shaped: \texttt{th} was \texttt{ZI}, so \texttt{I}$=$%
\texttt{h} gives \texttt{tHQt}$=$\texttt{th?t}, which resolves to \texttt{that} with
$\texttt{Q}\to\texttt{a}$. The word \texttt{OL} is \texttt{?s}-shaped after we guess the next
piece; \texttt{JXeLtOoF} is an eight-letter word ending the sentence after \texttt{the}, with
\texttt{e} in slot 3 and \texttt{t} in slot 5 -- \texttt{??e?t??n} -- which resolves to
\texttt{question} with $\texttt{J}\to\texttt{q}$, $\texttt{X}\to\texttt{u}$,
$\texttt{L}\to\texttt{s}$, $\texttt{O}\to\texttt{i}$ (and \texttt{F}$=$\texttt{n}, already
fixed). That also reads \texttt{OL}$=$\texttt{is}. Fill in everything:
\begin{quote}\ttfamily
to be or not to be that is the question
\end{quote}

\begin{quote}\ttfamily
TO BE OR NOT TO BE THAT IS THE QUESTION
\end{quote}

We never tried even one of the $4\times10^{26}$ keys. We recovered the substitution
letter by letter, using frequencies to seed guesses and word structure to confirm or refute
them. A few hundred characters of ciphertext is plenty; with a thousand, the attack is
mechanical.

\begin{whybox}
\textbf{Why the giant key space did not save the cipher.} Security is not the same as ``many
keys.'' The substitution cipher leaks a statistic -- the frequency of each letter -- that is
preserved exactly under the secret permutation. An attacker who can read off the answer never
has to search the key space at all. The modern restatement: a cipher is broken if the
ciphertext reveals \emph{any} information about the plaintext beyond its length. Counting
letters reveals an enormous amount. This is precisely the property Shannon would later forbid
with his definition of perfect secrecy, and which modern ciphers achieve by making the output
statistically indistinguishable from random bits.
\end{whybox}

\begin{implementbox}
A toy frequency counter. This sorts ciphertext letters by how often they appear, which is the
seed of the manual attack above.
\begin{lstlisting}[language=Python]
# LEARNING ONLY -- this is a teaching toy, never use it to protect anything.
from collections import Counter

ENGLISH = "etaoinshrdlcumwfgypbvkjxqz"  # roughly most-to-least common

def frequency_attack(ciphertext):
    letters = [c for c in ciphertext.upper() if c.isalpha()]
    counts = Counter(letters)
    # most common ciphertext letter -> guess it is 'E', and so on
    ranked = [c for c, _ in counts.most_common()]
    guess = {}
    for cipher_letter, plain_letter in zip(ranked, ENGLISH):
        guess[cipher_letter] = plain_letter
    decoded = "".join(guess.get(c, c.lower()) for c in ciphertext.upper())
    return decoded, guess

text = "ZG WT GK FGZ ZG WT ZIQZ OL ZIT JXTLZOGF"
print(frequency_attack(text))
# A first pass; you then hand-correct using word patterns, exactly as above.
\end{lstlisting}
Real attacks add bigram/trigram statistics and dictionary matching, but the core is this
histogram. Preserved letter frequencies are the fatal leak: any scheme that preserves letter proportions preserves the language's statistical signature.
\end{implementbox}

\begin{applicationbox}
\textbf{Where you still meet substitution.} The newspaper ``cryptoquote'' puzzle is exactly a
monoalphabetic substitution, solvable by the method above. More seriously, naive ``encryption''
that maps each byte to a fixed other byte -- some legacy file formats, some toy obfuscation in
malware, ROT13 in Usenet -- is a substitution cipher and falls to frequency analysis or simple
known-plaintext. Whenever a system claims to ``encrypt'' but a repeated input always produces
the same scrambled output, you are looking at a substitution-style leak. Modern block-cipher modes
(Chapter~13) exist specifically to eliminate this determinism.
\end{applicationbox}

\section{The Vigenere cipher: stacking shifts to flatten the fingerprint}
\label{ch02:vigenere}

If the problem with substitution is that each plaintext letter always maps to the same
ciphertext letter, the natural fix is to make the mapping \emph{change} as you go. The
Vigenere cipher does this with a repeating keyword. It was described by Giovan Battista
Bellaso in 1553, later misattributed to Blaise de Vigenere, and earned the nickname
\emph{le chiffre indechiffrable} -- the indecipherable cipher -- because it resisted frequency
analysis for three centuries.

\begin{focusbox}
\textbf{Three hundred years of false confidence.} From the mid-1500s until the mid-1800s the
Vigenere cipher was widely believed unbreakable. Charles Babbage broke it around 1854 but did
not publish (some think at the request of British intelligence during the Crimean War). The
Prussian infantry officer Friedrich Kasiski published the method independently in 1863, and it
bears his name. This chapter makes the point repeatedly: a cipher being unbroken for
a long time is not a proof that it is secure. It may only mean nobody clever has published yet.
\end{focusbox}

\subsection{How it works}

Pick a keyword, say \texttt{KEY}. Write it repeatedly under the plaintext, and shift each
plaintext letter by the corresponding key letter (treating \texttt{A}=0 as ``no shift'',
\texttt{B}=1, etc.). Formally, with key letters $k_0, k_1, \ldots, k_{\ell-1}$ of length
$\ell$,
\[
  c_i \equiv m_i + k_{(i \bmod \ell)} \pmod{26}.
\]
It is a Caesar cipher whose shift cycles through the keyword.

\begin{example}[Encrypting with \texttt{KEY}]
\label{ch02:vig-encrypt}
Encrypt \texttt{ATTACKATDAWN} with keyword \texttt{KEY} ($\texttt{K}=10,\texttt{E}=4,
\texttt{Y}=24$).
\[
\begin{array}{c|cccccccccccc}
\text{plain} & \texttt{A}&\texttt{T}&\texttt{T}&\texttt{A}&\texttt{C}&\texttt{K}&\texttt{A}&\texttt{T}&\texttt{D}&\texttt{A}&\texttt{W}&\texttt{N}\\
m            & 0&19&19&0&2&10&0&19&3&0&22&13\\
\text{key}   & \texttt{K}&\texttt{E}&\texttt{Y}&\texttt{K}&\texttt{E}&\texttt{Y}&\texttt{K}&\texttt{E}&\texttt{Y}&\texttt{K}&\texttt{E}&\texttt{Y}\\
k            & 10&4&24&10&4&24&10&4&24&10&4&24\\
m+k \bmod 26 & 10&23&17&10&6&8&10&23&1&10&0&11\\
\text{cipher}& \texttt{K}&\texttt{X}&\texttt{R}&\texttt{K}&\texttt{G}&\texttt{I}&\texttt{K}&\texttt{X}&\texttt{B}&\texttt{K}&\texttt{A}&\texttt{L}
\end{array}
\]
Result: \texttt{KXRKGIKXBKAL}.
\end{example}

Consider the three \texttt{T}'s in the plaintext (positions $2$, $3$, $8$). The
first became \texttt{X}, the second \texttt{R}, the third \texttt{X} again -- because the key
letter under each was \texttt{E}, \texttt{Y}, \texttt{E}. The same plaintext letter no longer
maps to a single ciphertext letter; it splits according to where the key cycle catches it.
(The four \texttt{A}'s, by contrast, all happen to fall on a key-\texttt{K} position and so all
become \texttt{K} -- a reminder that the flattening is statistical, not guaranteed letter by
letter.) The frequency fingerprint is blurred. Run a frequency count on a long Vigenere ciphertext and the histogram is far flatter
than English -- which is exactly why al-Kindi's attack stalls.

\begin{intuitionbox}
Vigenere is really $\ell$ independent Caesar ciphers interleaved: positions $0, \ell, 2\ell,
\ldots$ all use key letter $k_0$; positions $1, \ell+1, 2\ell+1, \ldots$ all use $k_1$; and so
on. If you knew where one of these ``columns'' lived, you would be back to a single Caesar
cipher, which you can crack in 25 tries. So the problem reduces to this: \emph{find the key length
$\ell$.} Once you know $\ell$, the cipher collapses into $\ell$ easy pieces.
\end{intuitionbox}

Figure~\ref{ch02:fig-vig-columns} makes the ``$\ell$ interleaved Caesars'' picture concrete on
the ciphertext \texttt{KXRKGIKXBKAL} from Example~\ref{ch02:vig-encrypt}, whose key
\texttt{KEY} has length $\ell=3$.

\begin{figure}[H]
\centering
\begin{tikzpicture}[font=\small,>=stealth]
  % Column headers
  \node[draw=sourceblue!80,fill=sourceblue!15,thick,rounded corners=2pt,minimum width=1.6cm,minimum height=0.7cm,align=center,font=\scriptsize\bfseries] (h0) at (0,0) {Col 0\\$k_0=\texttt{K}\,(+10)$};
  \node[draw=targetgreen!80,fill=targetgreen!15,thick,rounded corners=2pt,minimum width=1.6cm,minimum height=0.7cm,align=center,font=\scriptsize\bfseries] (h1) at (1.9,0) {Col 1\\$k_1=\texttt{E}\,(+4)$};
  \node[draw=quotientorange!80,fill=quotientorange!15,thick,rounded corners=2pt,minimum width=1.6cm,minimum height=0.7cm,align=center,font=\scriptsize\bfseries] (h2) at (3.8,0) {Col 2\\$k_2=\texttt{Y}\,(+24)$};

  % Ciphertext letters with indices
  % Col 0: pos 0, 3, 6, 9 -> K, K, K, K
  \foreach \r/\idx/\L in {1/0/K, 2/3/K, 3/6/K, 4/9/K} {
    \node[draw=sourceblue!50,fill=white,rounded corners=1pt,minimum width=1.4cm,minimum height=0.55cm,font=\bfseries] at (0,-\r*0.65) {\texttt{\L}\ \tiny[pos \idx]};
  }

  % Col 1: pos 1, 4, 7, 10 -> X, G, X, A
  \foreach \r/\idx/\L in {1/1/X, 2/4/G, 3/7/X, 4/10/A} {
    \node[draw=targetgreen!50,fill=white,rounded corners=1pt,minimum width=1.4cm,minimum height=0.55cm,font=\bfseries] at (1.9,-\r*0.65) {\texttt{\L}\ \tiny[pos \idx]};
  }

  % Col 2: pos 2, 5, 8, 11 -> R, I, B, L
  \foreach \r/\idx/\L in {1/2/R, 2/5/I, 3/8/B, 4/11/L} {
    \node[draw=quotientorange!50,fill=white,rounded corners=1pt,minimum width=1.4cm,minimum height=0.55cm,font=\bfseries] at (3.8,-\r*0.65) {\texttt{\L}\ \tiny[pos \idx]};
  }

  % Explanatory Card on Right
  \node[draw=structurepurple!50,fill=structurepurple!5,rounded corners=4pt,thick,align=left,font=\scriptsize,inner sep=7pt] (card) at (9.1,-1.45) {
    \textbf{Structural Decomposition:}\\
    -- Each vertical slice uses a single, fixed shift.\\
    -- Col 0: Pure Caesar shift with $k_0=10$.\\
    -- Col 1: Pure Caesar shift with $k_1=4$.\\
    -- Col 2: Pure Caesar shift with $k_2=24$.\\[4pt]
    \textbf{Complexity Collapse:}\\
    Instead of testing $26^3 = 17{,}576$ three-letter keys,\\
    an attacker tests $3 \times 25 = 75$ shifts independently.
  };

  \draw[->,thick,structurepurple] (4.8,-1.45) -- (card.west);
\end{tikzpicture}
\caption{A Vigenere ciphertext sliced into $\ell=3$ columns according to position index modulo $3$. Because every character in column $j$ is transformed by the exact same key letter $k_j$, finding the period $\ell$ decomposes the polyalphabetic cipher into $\ell$ independent, trivial monoalphabetic Caesar shifts.}
\label{ch02:fig-vig-columns}
\end{figure}

\subsection{Kasiski and Babbage: finding the key length}
\label{ch02:kasiski}

The Vigenere cipher hides the key, but it cannot hide a coincidence. Suppose a common word
like \texttt{THE} appears twice in the plaintext, and by luck both occurrences line up against
the \emph{same} part of the keyword. Then both encrypt to the \emph{same} ciphertext triple.
The gap between those two identical triples in the ciphertext must be a whole number of
keyword-lengths -- it is a multiple of $\ell$.

\begin{whybox}
\textbf{Why repeated patterns leak the key length.} If plaintext position $i$ and position $j$
hold the same letters \emph{and} are encrypted with the same key letters, the ciphertexts
match. Same key letters means $i \equiv j \pmod{\ell}$, i.e. the distance $j-i$ is a multiple
of $\ell$. So every repeated ciphertext chunk hands you a number that $\ell$ divides. Collect
several such distances; the key length is a common divisor of all of them -- usually their
greatest common divisor, or a small factor of it. This is the Kasiski test.
\end{whybox}

\begin{attackbox}
\textbf{Kasiski test, worked.} Suppose Eve, hunting for repeated trigrams in a long
ciphertext, finds the trigram \texttt{WMV} at positions $5$, $21$, and $53$ (counting from
$0$), and the trigram \texttt{GTH} at positions $9$ and $25$.

Distances between repeats:
\[
\begin{aligned}
\texttt{WMV}: &\quad 21-5 = 16, \quad 53-21 = 32, \quad 53-5 = 48,\\
\texttt{GTH}: &\quad 25-9 = 16.
\end{aligned}
\]
Now take the greatest common divisor of the collected distances:
\[
  \gcd(16, 32, 48, 16) = 16.
\]
The key length divides $16$, so $\ell \in \{1,2,4,8,16\}$. A length of $1$ is just a Caesar
cipher (ruled out, the histogram was flat), so the strong candidates are $\ell = 8$ or
$\ell = 16$. Eve tries the smallest plausible value first (Figure~\ref{ch02:fig-kasiski}).
\end{attackbox}

\begin{figure}[H]
\centering
\begin{tikzpicture}[font=\small,>=stealth]
  % Ciphertext timeline bar
  \draw[thick,gray!50,fill=gray!8,rounded corners=2pt] (0,0) rectangle (11.5,0.65);
  \node[left,font=\footnotesize\bfseries,gray!80!black] at (-0.15,0.325) {Ciphertext:};

  % WMV blocks at pos 5, 21, 53 (scaled to x = 1.0, 4.2, 9.8)
  \node[draw=sourceblue!85!black,fill=sourceblue!18,thick,rounded corners=1pt,minimum width=1.1cm,minimum height=0.58cm,font=\footnotesize\bfseries,text=sourceblue!95!black] (w1) at (1.1,0.325) {\texttt{WMV}};
  \node[above=1pt,font=\tiny\bfseries,sourceblue!90!black] at (w1.north) {pos $5$};

  \node[draw=sourceblue!85!black,fill=sourceblue!18,thick,rounded corners=1pt,minimum width=1.1cm,minimum height=0.58cm,font=\footnotesize\bfseries,text=sourceblue!95!black] (w2) at (4.3,0.325) {\texttt{WMV}};
  \node[above=1pt,font=\tiny\bfseries,sourceblue!90!black] at (w2.north) {pos $21$};

  \node[draw=sourceblue!85!black,fill=sourceblue!18,thick,rounded corners=1pt,minimum width=1.1cm,minimum height=0.58cm,font=\footnotesize\bfseries,text=sourceblue!95!black] (w3) at (9.9,0.325) {\texttt{WMV}};
  \node[above=1pt,font=\tiny\bfseries,sourceblue!90!black] at (w3.north) {pos $53$};

  % GTH blocks at pos 9, 25 (scaled to x = 2.5, 5.7)
  \node[draw=targetgreen!85!black,fill=targetgreen!18,thick,rounded corners=1pt,minimum width=1.1cm,minimum height=0.58cm,font=\footnotesize\bfseries,text=targetgreen!95!black] (g1) at (2.5,0.325) {\texttt{GTH}};
  \node[above=1pt,font=\tiny\bfseries,targetgreen!90!black] at (g1.north) {pos $9$};

  \node[draw=targetgreen!85!black,fill=targetgreen!18,thick,rounded corners=1pt,minimum width=1.1cm,minimum height=0.58cm,font=\footnotesize\bfseries,text=targetgreen!95!black] (g2) at (5.7,0.325) {\texttt{GTH}};
  \node[above=1pt,font=\tiny\bfseries,targetgreen!90!black] at (g2.north) {pos $25$};

  % Distance brackets / arrows
  \draw[<->,thick,sourceblue!85!black] (1.1,-0.35) -- node[below=1pt,font=\scriptsize\bfseries]{$\Delta_1 = 21 - 5 = 16$} (4.3,-0.35);
  \draw[<->,thick,sourceblue!85!black] (4.3,-0.35) -- node[below=1pt,font=\scriptsize\bfseries]{$\Delta_2 = 53 - 21 = 32$} (9.9,-0.35);
  \draw[<->,thick,targetgreen!85!black] (2.5,-1.05) -- node[below=1pt,font=\scriptsize\bfseries]{$\Delta_3 = 25 - 9 = 16$} (5.7,-1.05);

  % GCD conclusion box
  \node[draw=quotientorange!80!black,fill=quotientorange!12,rounded corners=3pt,thick,font=\footnotesize,inner sep=5pt] at (5.75,-1.95) {
    \textbf{Kasiski deduction:} $\gcd(16, 32, 16) = 16 \implies \text{keyword length } \ell \text{ divides } 16$
  };
\end{tikzpicture}
\caption{The Kasiski test. When identical plaintext trigrams happen to align with the same phase of the repeating keyword, they produce identical ciphertext trigrams (\texttt{WMV} and \texttt{GTH}). The spacing between repeats must be an integer multiple of the period $\ell$, so the greatest common divisor of those spacings exposes $\ell$.}
\label{ch02:fig-kasiski}
\end{figure}

Once Eve commits to a guess $\ell = 8$, she splits the ciphertext into $8$ columns by
position modulo $8$. Each column is a pure Caesar cipher, and she breaks each one in turn --
either by brute-forcing its $26$ shifts and reading for English, or, better, by frequency
analysis on that column alone. The most common letter in column $j$ is probably the encryption
of \texttt{E}, which pins down $k_j$.

\begin{example}[Solving one column]
\label{ch02:vig-column}
Suppose column $0$ (positions $0,8,16,\ldots$) has these letter counts, with \texttt{R} the
clear winner. Guess that \texttt{R} is the encryption of \texttt{E}. Then
\[
  k_0 \equiv \texttt{R} - \texttt{E} = 17 - 4 = 13 \pmod{26},
\]
so $k_0 = 13 = \texttt{N}$. The first keyword letter is \texttt{N}. Repeat the same
``most-common-letter-is-\texttt{E}'' deduction on the other seven columns and the whole keyword
falls out -- say \texttt{NIGHTOWL}. If a recovered letter looks wrong (the second-most-common
letter might really be \texttt{E}, with \texttt{T} the winner instead), Eve simply tries the
top two or three candidates per column and keeps the keyword that yields readable English.
\end{example}

\subsection{The Friedman test: counting coincidences}
\label{ch02:friedman}

There is a second, quantitative way to estimate the key length, due to William Friedman in
1920, that does not rely on spotting repeats by eye. It is worth seeing because it is the first
appearance of a statistic we will reuse constantly: the \emph{index of coincidence}.

\begin{definition}[Index of coincidence]
\label{ch02:ioc-def}
Given a text of length $n$ with letter counts $n_A, n_B, \ldots, n_Z$, the index of
coincidence is the probability that two letters drawn at random (without replacement) are
equal:
\[
  \mathrm{IC} \;=\; \frac{\sum_{x=A}^{Z} n_x (n_x - 1)}{n(n-1)}.
\]
\end{definition}

For text drawn uniformly at random over $26$ letters, $\mathrm{IC} \approx 1/26 \approx
0.0385$. For English, the lopsided frequencies push it up to about $\mathrm{IC} \approx
0.0667$. A monoalphabetic substitution \emph{preserves} the English value (it only renames
letters), so it sits near $0.066$ -- a dead giveaway. A Vigenere ciphertext with a long key
looks much closer to random, near $0.038$.

\begin{whybox}
\textbf{Why the index of coincidence finds the key length.} If you write the ciphertext into
$\ell$ columns and $\ell$ is the true key length, each column is a single Caesar shift of
English and therefore has English-like statistics, with $\mathrm{IC}\approx 0.066$. If your
guessed $\ell$ is wrong, the columns mix several different shifts and look random,
$\mathrm{IC}\approx 0.038$. So you sweep $\ell = 1, 2, 3, \ldots$, compute the average column
IC for each, and watch for the value that jumps up toward $0.066$. That $\ell$ is your key
length -- no eyeballing of repeated trigrams required. Friedman even gave a direct estimate:
\[
  \ell \;\approx\; \frac{0.027\, n}{(n-1)\,\mathrm{IC} - 0.038\, n + 0.066}.
\]
\end{whybox}

\begin{connectionbox}
The index of coincidence is a collision probability: the chance that two independently chosen
items land in the same bucket. That same quantity drives the birthday paradox and hash
collisions, covered in the companion probability chapters and revisited when we study hash
functions. The Friedman test asks ``how surprised am I by how many coincidences I
see?'' -- a question you answer with exactly the combinatorics used to bound a birthday attack.
\end{connectionbox}

\begin{implementbox}
Index of coincidence and a key-length sweep.
\begin{lstlisting}[language=Python]
# LEARNING ONLY -- toy code for understanding, not for protecting real data.
def index_of_coincidence(text):
    text = [c for c in text.upper() if c.isalpha()]
    n = len(text)
    if n < 2:
        return 0.0
    from collections import Counter
    counts = Counter(text)
    return sum(v * (v - 1) for v in counts.values()) / (n * (n - 1))

def guess_key_length(ciphertext, max_len=20):
    letters = [c for c in ciphertext.upper() if c.isalpha()]
    for ell in range(1, max_len + 1):
        columns = ["".join(letters[i::ell]) for i in range(ell)]
        avg_ic = sum(index_of_coincidence(col) for col in columns) / ell
        print(f"key length {ell:2d}: avg IC = {avg_ic:.4f}")
        # English-like (~0.066) flags the true length; random-like is ~0.038
\end{lstlisting}
Run this on a Vigenere ciphertext and the true key length is the row where the average index
of coincidence spikes toward $0.066$. Then break each column as a Caesar cipher.
\end{implementbox}

\begin{trapbox}
\textbf{``I'll just use a really long key.'' } Lengthening the Vigenere keyword raises the bar
but does not change the verdict. If the key is shorter than the message and repeats, Kasiski
and Friedman eventually find the period given enough ciphertext. The \emph{only} way a Vigenere
key is truly safe is if it is as long as the message, never repeats, and is random -- at which
point it stops being Vigenere and becomes the one-time pad (Chapter~9). The same lesson
applies throughout modern crypto: \textbf{a repeating keystream is a broken keystream.} This
holds whether the repeated value is a Vigenere keyword, a stream-cipher nonce, or an IV.
\end{trapbox}

\begin{applicationbox}
\textbf{The same bug, a century later.} The exact failure mode of a repeating Vigenere key
reappears in modern systems whenever a stream-cipher keystream gets reused. The 2001 break of
WiFi's WEP exploited IV reuse in RC4. The ``two-time pad'' mistake -- encrypting two messages
with the same keystream -- lets an attacker XOR the ciphertexts together and cancel the key,
leaving $m_1 \xor m_2$, which is then separated by the same language-statistics reasoning we
used on Vigenere columns. Microsoft's PPTP and several TLS missteps hit variants of this. If
you remember one thing from this section: never reuse a keystream, and treat any cipher whose
keystream repeats as already broken.
\end{applicationbox}

\section{Rotor machines and Enigma: mechanizing the substitution}
\label{ch02:enigma}

By the early twentieth century, cryptographers understood that a fixed substitution was fatal
and that you wanted the substitution to change with every letter. Doing that by hand with a
long keyword is slow and error-prone. The breakthrough was mechanical: build a machine whose
internal wiring \emph{is} a substitution, and make it step to a new substitution after every
keypress. These are \emph{rotor machines}, and the most famous is the German Enigma.

\begin{historybox}
\textbf{Enigma.} Patented by Arthur Scherbius around 1918 and adopted by the German military through the 1920s and 30s, Enigma was the primary cipher of the Wehrmacht, Luftwaffe, and Kriegsmarine in World War II. Its operators believed the $10^{23}$ rotor-and-plugboard configurations made it unbreakable. They were wrong: Polish mathematicians Marian Rejewski, Jerzy R\'o\.zycki, and Henryk Zygalski broke it in 1932, and after Poland shared its methods and replica machines with Britain and France in July 1939, Alan Turing, Gordon Welchman, and thousands of others at Bletchley Park industrialized the attack. Purpose-built machines (the Polish \emph{bomba}, the British \emph{Bombe}) turned cryptanalysis into an engineering discipline and shortened the war by years. Afterward, Britain kept Bletchley's work classified until the 1970s, obscuring Poland's breakthrough while selling captured rotor machines to former allies without revealing they could be read.
\end{historybox}

\subsection{What the machine does}

A rotor is a disk with $26$ contacts on each face, wired so that current entering at one
letter exits at some scrambled other letter -- a fixed substitution baked into copper. Stack
three rotors in series and you compose three substitutions. After each keypress, the rightmost
rotor turns one position, so the overall substitution changes every letter. When the right
rotor completes a full turn it kicks the middle rotor one step (like an odometer), and the
middle eventually kicks the left. Enigma added two more elements:

\begin{itemize}
  \item A \emph{reflector} at the far end that sent the current back through the rotors a
        second time. This made encryption and decryption the same operation (press the key,
        read the lamp) -- a great convenience.
  \item A \emph{plugboard} (\emph{Steckerbrett}) that swapped up to $10$ pairs of letters
        before and after the rotor stack, multiplying the key space enormously.
\end{itemize}

\begin{figure}[H]
\centering
\begin{tikzpicture}[font=\small,>=stealth,scale=0.92]
  % Components
  \node[draw=sourceblue!80,fill=sourceblue!12,thick,rounded corners=3pt,minimum width=1.5cm,minimum height=2.4cm,align=center] (kb) at (0,0) {\textbf{Keyboard}\\$\&$ \textbf{Lamps}\\[4pt]\tiny Press: \texttt{A}\\\tiny Light: \texttt{W}};
  \node[draw=structurepurple!80,fill=structurepurple!12,thick,rounded corners=3pt,minimum width=1.5cm,minimum height=2.4cm,align=center] (plug) at (2.4,0) {\textbf{Plugboard}\\(\emph{Stecker})\\[4pt]\tiny $\texttt{A} \leftrightarrow \texttt{J}$\\\tiny $\texttt{P} \leftrightarrow \texttt{W}$};
  
  \node[draw=targetgreen!80,fill=targetgreen!12,thick,rounded corners=3pt,minimum width=1.35cm,minimum height=2.4cm,align=center] (r3) at (4.6,0) {\textbf{Rotor 3}\\(Fast)\\[4pt]\tiny Steps every\\\tiny keypress};
  \node[draw=targetgreen!80,fill=targetgreen!12,thick,rounded corners=3pt,minimum width=1.35cm,minimum height=2.4cm,align=center] (r2) at (6.6,0) {\textbf{Rotor 2}\\(Middle)\\[4pt]\tiny Turnover\\\tiny ring notch};
  \node[draw=targetgreen!80,fill=targetgreen!12,thick,rounded corners=3pt,minimum width=1.35cm,minimum height=2.4cm,align=center] (r1) at (8.6,0) {\textbf{Rotor 1}\\(Slow)\\[4pt]\tiny Turnover\\\tiny ring notch};
  
  \node[draw=quotientorange!80,fill=quotientorange!15,thick,rounded corners=3pt,minimum width=1.8cm,minimum height=2.4cm,align=center] (ref) at (10.9,0) {};
  \node[align=center] at (11.0,0) {\textbf{Reflector}\\(\emph{UKW})\\[4pt]\tiny Pairwise\\\tiny loopbacks};

  % Forward path (Upper lane, blue)
  \draw[->,thick,sourceblue!90!black] ($(kb.east)+(0,0.65)$) -- node[above,font=\tiny\bfseries]{$\texttt{A}$} ($(plug.west)+(0,0.65)$);
  \draw[->,thick,sourceblue!90!black] ($(plug.east)+(0,0.65)$) -- node[above,font=\tiny\bfseries]{$\texttt{J}$} ($(r3.west)+(0,0.65)$);
  \draw[->,thick,sourceblue!90!black] ($(r3.east)+(0,0.65)$) -- ($(r2.west)+(0,0.65)$);
  \draw[->,thick,sourceblue!90!black] ($(r2.east)+(0,0.65)$) -- ($(r1.west)+(0,0.65)$);
  \draw[->,thick,sourceblue!90!black] ($(r1.east)+(0,0.65)$) -- ($(ref.west)+(0,0.65)$);

  % U-turn at reflector left edge (clear of text)
  \draw[thick,quotientorange!90!black] ($(ref.west)+(0,0.65)$) -- ($(ref.west)+(0.10,0.65)$) arc (90:-90:0.12cm and 0.65cm) -- ($(ref.west)+(0,-0.65)$);

  % Return path (Lower lane, orange dashed)
  \draw[<-,thick,dashed,quotientorange!90!black] ($(r1.east)+(0,-0.65)$) -- ($(ref.west)+(0,-0.65)$);
  \draw[<-,thick,dashed,quotientorange!90!black] ($(r2.east)+(0,-0.65)$) -- ($(r1.west)+(0,-0.65)$);
  \draw[<-,thick,dashed,quotientorange!90!black] ($(r3.east)+(0,-0.65)$) -- ($(r2.west)+(0,-0.65)$);
  \draw[<-,thick,dashed,quotientorange!90!black] ($(plug.east)+(0,-0.65)$) -- node[below,font=\tiny\bfseries]{$\texttt{P}$} ($(r3.west)+(0,-0.65)$);
  \draw[<-,thick,dashed,quotientorange!90!black] ($(kb.east)+(0,-0.65)$) -- node[below,font=\tiny\bfseries]{$\texttt{W}$} ($(plug.west)+(0,-0.65)$);

  % Explanatory legend
  \node[font=\tiny,sourceblue!90!black,above] at (5.6,1.4) {\textbf{Forward current:} Key \texttt{A} $\to$ Plugboard swap to \texttt{J} $\to$ Through Rotors 3, 2, 1};
  \node[font=\tiny,quotientorange!90!black,below] at (5.6,-1.4) {\textbf{Reflected return current:} Through Rotors 1, 2, 3 $\to$ Plugboard swap $\texttt{P}\to\texttt{W}$ $\to$ Lights lamp \texttt{W}};
\end{tikzpicture}
\caption{The Enigma electro-mechanical signal path. Current leaves the pressed key, traverses the plugboard and three rotors, bounces off the reflector, and returns back through the rotors along a completely distinct electrical circuit to illuminate the lamp. Because the reflector connects contacts only in disjoint pairs, no electrical signal can ever return to the key that initiated it.}
\label{ch02:fig-enigma-wiring}
\end{figure}

The daily key was: which rotors in which order, their starting positions, the ring settings,
and the plugboard pairings. Counting it all up, the Enigma key space is around $10^{23}$ --
roughly $77$ bits. By the brute-force logic of Section~\ref{ch02:caesar}, that should be safe.
It was not, and the reasons show how real systems fail.

\subsection{The flaws that sank it}
\label{ch02:enigma-flaws}

\begin{attackbox}
\textbf{No letter ever encrypts to itself.} The reflector wiring guarantees that the output
letter is always different from the input letter. This sounds harmless, but it was a serious weakness. It
let codebreakers \emph{slide a guessed plaintext (a ``crib'') along the ciphertext} and instantly
discard any alignment where the crib and ciphertext shared a letter in the same position. A
single self-mapping anywhere kills that alignment. With common cribs -- weather reports always
contained \texttt{WETTER}, messages often ended \texttt{HEILHITLER} -- this pruned the search
enormously before a single rotor was tried.
\end{attackbox}

\begin{figure}[H]
\centering
\begin{tikzpicture}[font=\small,scale=0.92]
  % Intercepted Ciphertext strip
  \node[font=\bfseries\footnotesize,sourceblue,anchor=east] at (-0.4,1.8) {Ciphertext $C$:};
  \foreach \i/\L in {0/E, 1/W, 2/K, 3/Z, 4/L, 5/P, 6/A, 7/Q, 8/M, 9/X, 10/N} {
    \node[draw=sourceblue!70,fill=sourceblue!10,thick,rounded corners=1pt,minimum width=0.7cm,minimum height=0.6cm,font=\bfseries] (c\i) at (\i*0.8,1.8) {\texttt{\L}};
    \node[above,font=\tiny,gray] at (c\i.north) {\i};
  }

  % Trial 1 (Offset 1) - Collision at pos 1! Placed immediately below C so vertical clash arrow is unobstructed
  \node[font=\bfseries\tiny,trapred!90!black,anchor=east] at (-0.4,0.65) {Offset 1:};
  \foreach \i/\L in {1/W, 2/E, 3/T, 4/T, 5/E, 6/R} {
    \ifnum\i=1
      \node[draw=trapred,fill=trapred!25,very thick,rounded corners=1pt,minimum width=0.7cm,minimum height=0.55cm,font=\bfseries,text=trapred!90!black] (t1\i) at (\i*0.8,0.65) {\texttt{\L}};
    \else
      \node[draw=gray!50,fill=gray!10,rounded corners=1pt,minimum width=0.7cm,minimum height=0.55cm,font=\bfseries] (t1\i) at (\i*0.8,0.65) {\texttt{\L}};
    \fi
  }
  % Clash highlight line
  \draw[<->,very thick,trapred] (c1.south) -- (t11.north);
  \node[draw=trapred!80,fill=trapred!15,rounded corners=2pt,font=\tiny\bfseries,right=6pt] at (t16.east) {$\times$\ Rejected: \texttt{W} matches \texttt{W} at pos 1 (impossible)};

  % Trial 2 (Offset 0)
  \node[font=\bfseries\tiny,targetgreen!80!black,anchor=east] at (-0.4,-0.25) {Offset 0:};
  \foreach \i/\L in {0/W, 1/E, 2/T, 3/T, 4/E, 5/R} {
    \node[draw=targetgreen!70,fill=targetgreen!10,rounded corners=1pt,minimum width=0.7cm,minimum height=0.55cm,font=\bfseries] (t0\i) at (\i*0.8,-0.25) {\texttt{\L}};
  }
  \node[draw=targetgreen!80,fill=targetgreen!15,rounded corners=2pt,font=\tiny\bfseries,right=6pt] at (t05.east) {\checkmark\ Candidate: zero collisions};

  % Trial 3 (Offset 2)
  \node[font=\bfseries\tiny,targetgreen!80!black,anchor=east] at (-0.4,-1.1) {Offset 2:};
  \foreach \i/\L in {2/W, 3/E, 4/T, 5/T, 6/E, 7/R} {
    \node[draw=targetgreen!70,fill=targetgreen!10,rounded corners=1pt,minimum width=0.7cm,minimum height=0.55cm,font=\bfseries] (t2\i) at (\i*0.8,-1.1) {\texttt{\L}};
  }
  \node[draw=targetgreen!80,fill=targetgreen!15,rounded corners=2pt,font=\tiny\bfseries,right=6pt] at (t27.east) {\checkmark\ Candidate: zero collisions};
\end{tikzpicture}
\caption{Crib-dragging against Enigma ciphertext. Because Enigma's electrical reflector prevents any letter from ever encrypting to itself, testing a suspected crib (such as \texttt{WETTER}) against the intercepted stream is instantaneous: any offset that aligns identical letters (like the \texttt{W} collision at offset 1) is ruled out immediately with zero rotor testing.}
\label{ch02:fig-enigma-crib}
\end{figure}

\begin{attackbox}
\textbf{Operator habits and key reuse.} Procedure required each message to carry a per-message
rotor start, encrypted under the day key. Early procedure had the operator type this three-letter
\emph{message key twice} (e.g. \texttt{QWE QWE}) so the receiver could check it. Typing it twice
created a fixed relationship between ciphertext positions $1$\&$4$, $2$\&$5$, $3$\&$6$ -- exactly
the kind of repeated structure we exploited in Vigenere. Rejewski turned this into a system of
permutation equations and solved for the wiring. Beyond that, operators reused favorite settings
(\texttt{CIL} for a girlfriend's name, \texttt{BER} for Berlin), chose non-random plugboard
pairings, and sent stereotyped messages. Each of these habits gave the codebreakers something to exploit.
\end{attackbox}

\begin{whybox}
\textbf{Why ``no self-encryption'' is a leak, in our language.} A secure cipher's output should
look like a uniformly random string, independent of the plaintext except for length. Enigma's
ciphertext was \emph{not} uniform: the event ``ciphertext letter equals plaintext letter'' had
probability exactly $0$, whereas a random cipher would show it about $1/26$ of the time. That is
information about the plaintext leaking into the ciphertext -- forbidden by the security
definitions we build in later chapters. The Bombe was, in effect, a machine for exploiting that
single missing $1/26$, chaining crib-driven deductions through the rotors and using
contradictions (a letter forced to two different plug partners) to reject wrong rotor settings
at electrical speed.
\end{whybox}

\begin{intuitionbox}
The Bombe did not try all $10^{23}$ keys. Given a crib -- a guessed plaintext aligned under
ciphertext -- it set up a \emph{logical chain}: if plug-partner of \texttt{A} is \texttt{T},
then by following the rotor wiring this forces the partner of \texttt{T} to be \texttt{G},
which forces \ldots and if the chain ever loops back and contradicts itself, that rotor setting
is impossible. The Bombe mechanized this deduction across all rotor positions, stopping only at
settings with no contradiction. The plugboard, which had made the key space huge, barely
mattered once a good crib was in hand -- because the deductions \emph{ignored} the plug swaps
until the very end. A huge key space bought the Germans little against an attack that never
searched it.
\end{intuitionbox}

\begin{focusbox}
\textbf{The Polish head start.} It is easy to picture Enigma as a Bletchley story, but the
decisive early work was Polish. Marian Rejewski, handed the doubled-message-key flaw and some
intelligence on the wiring, reconstructed Enigma's internal connections in 1932 using group
theory -- specifically, the cycle structure of permutations, which is invariant under the
unknown plugboard and so let him peel the plugboard away from the rotor problem. Poland built
electromechanical \emph{bomba} machines and read German traffic for years. In July 1939, weeks
before the invasion, Poland handed its results and replica machines to France and Britain. The
British work built on that foundation. The lesson for a working cryptographer:
the math that breaks your system may already exist in someone's unpublished work.
\end{focusbox}

\begin{applicationbox}
\textbf{Why Enigma still matters to you.} Enigma's downfall was not a weak core algorithm in
isolation -- it was the \emph{system}: stereotyped messages giving cribs, a structural quirk
(no self-encryption), key-management procedures that reused and exposed message keys, and human
operators who chose predictable settings. Every modern breach you will read about follows the
same template. Heartbleed, the WEP break, the Sony PS3 ECDSA nonce reuse, countless
implementation CVEs -- the cipher math was usually fine; the system around it leaked. When you
deploy crypto, your threat surface is the protocol, the randomness source, the key lifecycle,
and the human, not just the cipher. Enigma shows clearly that you cannot reason about a
cipher without reasoning about how people use it.
\end{applicationbox}

\subsection{A tiny rotor by hand}
\label{ch02:toy-rotor}

To see the mechanism, build a one-rotor toy on a $5$-letter alphabet \texttt{A B C D E}
$=0,1,2,3,4$. Let the rotor be the fixed substitution
\[
  S: \quad \texttt{A}\to\texttt{C},\ \texttt{B}\to\texttt{E},\ \texttt{C}\to\texttt{A},\
  \texttt{D}\to\texttt{B},\ \texttt{E}\to\texttt{D},
\]
and let the rotor ``step'' by adding $1$ (mod $5$) to the input position before applying $S$
and subtracting it after -- a stripped-down version of how a rotor's rotation shifts its whole
map. Encrypting letter $m_i$ at step $i$ (starting $i=0$):
\[
  c_i = \big(S\big((m_i + i)\bmod 5\big) - i\big)\bmod 5,
\]
where we read $S$ as a function on $\{0,1,2,3,4\}$: $S(0)=2,S(1)=4,S(2)=0,S(3)=1,S(4)=3$.

Encrypt \texttt{BAD} $=(1,0,3)$:
\[
\begin{aligned}
i=0:&\ \ S((1+0)\bmod5)-0 = S(1)=4 = \texttt{E},\\
i=1:&\ \ S((0+1)\bmod5)-1 = S(1)-1 = 4-1 = 3 = \texttt{D},\\
i=2:&\ \ S((3+2)\bmod5)-2 = S(0)-2 = 2-2 = 0 = \texttt{A}.
\end{aligned}
\]
So \texttt{BAD} $\to$ \texttt{EDA}. The two different plaintext letters here all pass through
the \emph{same} wiring $S$, yet because the rotor stepped, identical inputs at different
positions would land differently -- the property hand-cranked Vigenere also had, now produced
by a machine. Watch, too, that with this $S$ a letter \emph{can} map to itself at some
positions; the real Enigma's reflector is what forced the never-itself property that doomed it.

\begin{implementbox}
The toy rotor, in code, so you can watch the substitution change at each step.
\begin{lstlisting}[language=Python]
# LEARNING ONLY -- a 5-letter cartoon of a rotor, not a real cipher.
S = [2, 4, 0, 1, 3]          # the fixed wiring (a permutation of 0..4)
S_inv = [S.index(j) for j in range(5)]

def enc(plain):
    out = []
    for i, ch in enumerate(plain):
        m = ord(ch) - ord('A')
        c = (S[(m + i) % 5] - i) % 5
        out.append(chr(c + ord('A')))
    return "".join(out)

def dec(cipher):
    out = []
    for i, ch in enumerate(cipher):
        c = ord(ch) - ord('A')
        m = (S_inv[(c + i) % 5] - i) % 5
        out.append(chr(m + ord('A')))
    return "".join(out)

print(enc("BAD"))        # -> EDA
print(dec(enc("BAD")))   # -> BAD
\end{lstlisting}
Add more rotors (compose several permutations, each stepping at a different rate) and a
reflector, and you have the skeleton of Enigma. Add a plugboard (a fixed swap applied before
and after) and you have its shape. None of this is secure -- it is here so the machine stops
being a black box.
\end{implementbox}

\section{What three thousand years of failure teaches}
\label{ch02:lessons}

Taken together, these ciphers show a clear pattern. Every cipher in this chapter
was built by capable people, trusted by serious institutions, and broken by attackers who
did not behave as the designers expected. The recurring lessons are worth stating
plainly, because everything that follows builds on them.

\begin{focusbox}
\textbf{Five lessons that outlived their ciphers.}
\begin{enumerate}
  \item \emph{The key space must be too large to search} -- but a large key space guarantees
        nothing on its own (Section~\ref{ch02:substitution}, Section~\ref{ch02:enigma}).
        Caesar lost to brute force; substitution and Enigma had enormous key spaces and lost
        anyway, because the attacker never searched.
  \item \emph{Ciphertext must not leak statistics about plaintext.} Letter frequencies sank
        substitution; the missing $1/26$ self-map sank Enigma; a flatter-but-not-flat histogram
        and repeated patterns sank Vigenere. Modern security definitions formalize exactly this:
        ciphertext should be indistinguishable from random.
  \item \emph{Reuse is the enemy.} A repeating Vigenere keyword, a doubled Enigma message key,
        a reused keystream -- all the same mistake. Keys and keystreams must not repeat.
  \item \emph{Kerckhoffs was right: assume the enemy knows the system.} Every cipher here was
        eventually understood by the attacker; only key secrecy ever mattered, and key
        \emph{management} -- generation, distribution, rotation, operator discipline -- was
        where systems actually died.
  \item \emph{Amateurs lose.} ``It looks scrambled and nobody has broken it'' is not security.
        Vigenere went three centuries unbroken and was still broken. You earn confidence from
        analysis and proof, not from the absence of a published attack.
\end{enumerate}
\end{focusbox}

\begin{aestheticbox}
This history also shows a broader change in method. The early ciphers were defended by
\emph{obscurity and effort} -- surely no one will try all the shifts, surely no one will count
the letters, surely no one will notice the doubled key. Each defense was an appeal to the
attacker's laziness or ignorance. The modern field replaces that with something sturdier: we
state precisely what the attacker can do, we define exactly what ``broken'' means, and we prove
-- or reduce to a hard mathematical problem -- that breaking the cipher is as hard as solving
that problem. The rest of this book builds that foundation. Still, keep
the lessons of Caesar, Vigenere, and Enigma in mind: every one of them looked ``visibly
fine'' to its builders.
\end{aestheticbox}

\begin{connectionbox}
The thread from here forward: substitution's frequency leak motivates Shannon's \emph{perfect secrecy} and the one-time pad (Chapter~9); Vigen\`ere's repeating keystream motivates \emph{stream ciphers} and the absolute rule against keystream reuse (Chapter~10); Enigma's mechanized, ever-changing substitution is the ancestor of the \emph{block cipher} and its round structure (Chapter~11); and the recurring ``large key space, broken anyway'' is why we define security against polynomial-time adversaries rather than counting keys. Each modern construction is, in part, a direct answer to one of the failures in this chapter.

Every cipher here also had something in common: Caesar shifted letters around a ring of $26$ positions ($c \equiv m + k \pmod{26}$), Vigen\`ere cycled through shifts, and Enigma stepped contacts around a $26$-letter dial. In Chapter~3 we put this clockwork arithmetic on solid mathematical footing -- congruence classes, modular multiplication, and modular inverses -- the arithmetic that every modern cipher uses.
\end{connectionbox}

\section*{Exercises}
\addcontentsline{toc}{section}{Exercises}

\begin{exercise}[Caesar by hand]
\label{ch02:ex-caesar}
Encrypt the plaintext \texttt{CRYPTO} with a Caesar shift of $k=7$. Then decrypt the
ciphertext \texttt{ROVVY} (also $k=7$) back to plaintext. Show the numeric steps modulo $26$.
\end{exercise}

\begin{exercise}[Brute-forcing Caesar]
\label{ch02:ex-brute}
You intercept \texttt{IBWFGVO}. Knowing only that it is a Caesar cipher of an English word,
find the key and the plaintext by trying shifts. How many shifts did you need to test in the
worst case, and why is $26$ a fatal key-space size?
\end{exercise}

\begin{exercise}[Frequency attack]
\label{ch02:ex-freq}
The following is a monoalphabetic substitution of an English sentence:
\texttt{ZIT DTF QKT QSS OF ZIT ZTFZ}. Count the letters, identify the likely encryption
of \texttt{E} and \texttt{T}, use the short words and the repeated three-letter word
\texttt{ZIT} and the doubled-letter word \texttt{QSS}, and recover the plaintext. State your
full substitution map.
\end{exercise}

\begin{exercise}[Why $26!$ does not help]
\label{ch02:ex-keyspace}
The monoalphabetic key space has size $26! \approx 4\times10^{26}$, about $88$ bits. Explain in
two or three sentences why this enormous number provides essentially no security, and connect
your explanation to the notion that ciphertext must not leak statistics about plaintext.
\end{exercise}

\begin{exercise}[Vigenere encryption]
\label{ch02:ex-vig-enc}
Encrypt \texttt{MEETATDAWN} with the Vigenere keyword \texttt{LEMON}. Write the key letters
under the plaintext, give the numeric shifts, and produce the ciphertext. Then verify by
decrypting your own output.
\end{exercise}

\begin{exercise}[Kasiski test]
\label{ch02:ex-kasiski}
In a Vigenere ciphertext you find the trigram \texttt{XQR} repeated at positions $4$, $22$, and
$40$, and the trigram \texttt{PLM} repeated at positions $7$ and $31$. Compute all pairwise
distances, take their greatest common divisor, and list the candidate key lengths. Which would
you try first, and why?
\end{exercise}

\begin{exercise}[Index of coincidence]
\label{ch02:ex-ioc}
A ciphertext of length $n=200$ has letter counts that yield $\sum_x n_x(n_x-1) = 2600$. Compute
the index of coincidence. Is this text closer to English ($\approx 0.066$) or to random
($\approx 0.038$)? What does that suggest about whether it is a monoalphabetic substitution or a
Vigenere cipher with a long key?
\end{exercise}

\begin{exercise}[Two-time pad / keystream reuse]
\label{ch02:ex-twotime}
Alice encrypts two English messages $m_1$ and $m_2$ with the \emph{same} Vigenere keyword (or
equivalently the same keystream), producing $c_1$ and $c_2$. Show that $c_1 - c_2 \equiv m_1 -
m_2 \pmod{26}$ position by position, so the key cancels. Explain how an attacker recovers
$m_1$ and $m_2$ from $m_1 - m_2$ using language statistics, and name a modern system that
suffered this exact failure.
\end{exercise}

\begin{exercise}[Enigma's fatal quirk]
\label{ch02:ex-enigma-self}
The Enigma reflector guaranteed that no letter ever encrypts to itself. Suppose you suspect a
ciphertext segment \texttt{LXGPQ} is the encryption of the crib \texttt{ANGRI} (a fragment of
\texttt{ANGRIFF}). Check whether this alignment is possible under Enigma, and if not, explain
how this single rule lets a codebreaker discard crib positions in bulk before testing any rotor
setting.
\end{exercise}

\begin{exercise}[Toy rotor]
\label{ch02:ex-rotor}
Using the five-letter toy rotor of Section~\ref{ch02:toy-rotor} with
$S(0)=2,S(1)=4,S(2)=0,S(3)=1,S(4)=3$ and the stepping rule $c_i = (S((m_i+i)\bmod 5)-i)\bmod 5$,
encrypt the plaintext \texttt{CAB}. Then decrypt your result with the inverse rule to confirm
you recover \texttt{CAB}.
\end{exercise}

\begin{exercise}[Design a defense]
\label{ch02:ex-design}
For each of the three breaks in this chapter -- frequency analysis on substitution, the Kasiski
test on Vigenere, and crib-dragging on Enigma -- name the single property of the cipher that
made the attack possible, and state the modern design principle (from the five lessons in
Section~\ref{ch02:lessons}) that eliminates it. Keep each answer to one or two sentences.
\end{exercise}

\begin{exercise}[Prove a small fact about shifts]
\label{ch02:ex-prove}
Show that composing two Caesar ciphers, first a shift by $a$ then a shift by $b$, is the same as
a single Caesar shift by $(a+b)\bmod 26$. Conclude that the set of Caesar shifts forms a group
under composition, and explain why this means ``double Caesar'' adds no security at all.
\end{exercise}
